\documentclass[draft,landscape, 11pt, oneside]{report}
\usepackage[a4paper, left=2cm, right=2cm, top=2.5cm]{geometry}
\usepackage[T1]{fontenc}
\usepackage{lmodern}
\usepackage{hyperref}
\usepackage{titlesec}
\usepackage{scrextend}
\usepackage{array}
\usepackage{fontspec}
\usepackage[inline]{showlabels}
\setmainfont{FiraSans-SemiBold}
\setlength{\parindent}{0pt}
\setlength{\parskip}{1em}
\setlength{\leftmargin}{-4cm}
\setlength{\paperheight}{210mm}
\setlength{\paperwidth}{297mm}
\setlength{\textwidth}{26cm}
\setlength{\footskip}{2.5cm}
\renewcommand{\familydefault}{\sfdefault}

\newenvironment{flowtext}{\addmargin[0cm]{7cm}}{\endaddmargin} % standard text width (reduced for layout)

\title{[DRAFT] User Manual for the GliGli upgraded Prophet 600}
\author{Florian Merz (Editor)}
\pagestyle{myheadings}\markright{[DRAFT] User Manual for the GliGli upgraded Prophet 600}
\makeatletter
\renewcommand\chapter{\pagestyle{myheadings}\markright{[DRAFT] User Manual for the GliGli upgraded Prophet 600}\global\@topnum\z@\@afterindentfalse\secdef\@chapter\@schapter}
\makeatother
\titleformat{\chapter}[display]{\pagebreak\LARGE\bfseries}{}{0.0cm}{}
\titlespacing{\chapter}{0pt}{*0}{*0}
%\titleformat{\section}[display]{\Large\bfseries}{}{0.0cm}{}
\titlespacing{\section}{0pt}{*0}{*0}

\begin{document}

%\maketitle
\tableofcontents


\pagebreak
\chapter{Introduction to the GliGli Prophet 600 Firmware Upgrade}

\begin{flowtext}

\section{What you can expect from the Firmware Upgrade}

The credit of the original idea and implementation of the Prophet 600 firmware upgrade belongs to GliGli\cite{gligli}. He describes his motivation as follows:

\begin{addmargin}[2cm]{1cm}

\textit{I love vintage analog synthesizers, and to be honest, my dream synth would be a Prophet 5, but when I heard what the Prophet 600 was capable of tone wise, I immediately thought its major weaknesses -- the lousy computer part, software envelopes and LFO -- could become its strength with a remake; basically the whole internal synth in voltage-controlled from a nice 14bit DAC, so with a fast modern micro controller, it could become awesome, maybe even better than a Prophet 5!}

\end{addmargin}

The original Sequential Circuits Prophet 600 is run by a Z80 processor. This processor and firmware written for it reach their limits when computing the voltages, scanning user inputs and managing MIDI events making the instrment glitchy and non-responsive. GliGli reversed engineered the original firmware for the Z80 and found that the Z80 can be replaced by a (more) modern Teensy 2.0++ USB programmable microprocessor \cite{teensy}. The Teensy not only fits in to the 2x20 pin standard header but could also be programmed to take the role of Z80 with only slight modifications. In this way the limitations of the original Prophet 600 could be overcome - but more than that, it opened the possibility for new features (which may or may not have been on the feature list of Sequential Circuits). There are various modifications and modern CPU "implants" for some vintage analog synthesizers. One decisive advantage of the Teensy approach to the Prophet 600 is that it is an easy-to-install and non-destructive firmware drop-in replacement\footnote{There are two versions of the Prophet 600. One has the Z80 on a socket. In this case replacement is truly "drop-in". The other version has the Z80 soldered to the circuit board. In this case the Z80 needs to be removed a socket has to be soldered in its place first.}. So it is safe to try it out because one can always remove the Teensy board and restore the Z80 without technical expertise or training.

Developing the software was a major effort with low level functions at the core which deal with setting voltages to the electronic components, processing clock alignment, MIDI implementation, decoding the user inputs from dials, buttons and the keyboard and last but not least the "musical" coding of envelopes, modulations and sequencing functions. Many people have contributed, many out of curiosity and passion for synthesizers, most out of excitement that such a project could not only be conjured up but also be implemented and brought to life. The hardware of the Prophet 600 certainly justifies the effort: other than many DCO synths of the time the Prophet 600 is a 6 voice analog synthesizer driven by 12 voltage controlled oscillators (VCOs) and 24db low pass filters, all implemented using the famous Curtis chips\cite{curtis}. Furthermore, the hardware provides hard sync, ring modulation and modulation of cut-off by VCO. It is very versatile.

The following is a compact summary of the improved and the  new features which the upgraded Prophet 600 provides on this hardware. The manual refers to version 2.2 and differences between versions are marked or commented in the text. 

\begin{itemize}
  \setlength\itemsep{0cm}
  \item{General features}
  \subitem Greater resolution of many of the sound parameters with an improved refresh rate that is making the instrument much more responsive
  \subitem Faster and smoother ADSR envelopes with support two speed regimes (slow, fast) and two shapes (linear, exponential)
  \subitem positive and negative envelope amount settings for filter and poly-mod
  \subitem A new LFO function generator with a wider frequency range from one cycle every 20 seconds to about 60Hz
  \subitem Assignable, Random And Up/Down Arpeggiator sync-able to MIDI clock
  \subitem Multiple note assignment modes including last/low/high note priority
  \subitem Full MIDI In control including amplitude and filter velocity sensitivity with an external keyboard controller, continuous Controllers (CC) of all sound parameters, program change (PC) to choose current preset
  \subitem A dedicated vibrato which can be controlled by the modulation wheel, or can start progressively after a fixed amount of time
  \subitem Keyboard mode chord in addition to polyphonic and unison  
  \subitem Four new waveforms in addition to the original triangle and square including sine, random stepped, noise (like on the original Prophet 5, but non-periodic) and sawtooth (ramp up)
  \subitem Mix Overdrive which now allows the output from both oscillators to drive the mix VCAs A and B harder as well as the Curtis 4 pole filter resulting in new sonic possibilities
  \subitem Pitch Wheel interval selection of plus/minus one octave, a whole tone, a minor third and a fifth
  \subitem Pitch Wheel reassignment to the VCF and Volume or off   
  \subitem Modulation wheel intensity setting from Maximum to Medium to Minimum
  \subitem A new and improved tuning procedure
  \subitem Octave, chromatic and free Oscillator course pitch control
  \subitem Debounce feature that prevents unintended re-triggering caused by the old keyboard
  \item Features with version 2.1 RC3
  \subitem Polyphonic step sequencer
  \subitem Per note tuning
  \item Features with version 2.2
  \item Flexible envelope routing allowing the poly-mod to by modulated by either envelope generator
  \item Possibility to sync the LFO to the thythm of the arpeggiator or sequencer
\end{itemize}

\section{Versions of the firmware upgrade and compatibility}

The Firmware Upgrade project has been set up by GliGli as an open source project with a public code repository on GitHub\cite{repository}. The first stable version of the Firmware Upgrade was version 2.0 \cite{versiontwo}. An intermediate release (or at least something close to a release) version 2.1 Release Candidate 3 has been available from September 5th 2015 on \cite{gligli}. The current release is version 2.2. 

Upgrading the firmware can be done by SysEx (or directly using the USB port of the Teensy board), see section \ref{fwupgrade}. When upgrading the Prophet 600 patch data and settings are preserved. Note that since this is an open source project different variants of the firmware upgrade have been developed by users. In the main development branch, however, backward compatibility is intended to be ensured. This applies in particular to patches and settings as stored in the Prophet 600 memory or dumped SysEx files. It also applies to the panel overlay which some users have acquired to reflect the (slightly) altered function of some dials. 

Despite the fact that many have contributed and have spent considerable time to make this software as useful and consistent as possible and despite the fact that many users continue to enjoy the upgrade, the usage of the software and this documentation are understood to be at the user's own risk.

\section{Credits and license}

(this needs to be done)

\end{flowtext}

\chapter{Finding your way around your new  Prophet 600}    

\begin{flowtext}

\section{Display}

The 2 digit 7 segment display is used for different purposes.

\begin{itemize}
  \setlength\itemsep{0cm}
  \item Numeric controls: if the synth is in live mode (see section \ref{uimode}) and the user turns a dial, the display shows the value of the dial. A decimal point between the two digits indicates a negative value. Values range from -50 to 48 (signed parameters) or 0 to 98 (positive parameters). Note that this is a limitation of the display only and that the true resolution of the underlying parameters is higher for many parameters.  
  \item Parameter names and values: additional parameters are "hidden" in menus accessible via the key pad. The chosen parameter name and value are scrolled through the 2 digit display
  \item Messages such as the welcome message (software version) and confirmative messages (e.g. note of completion after a SysEx dump) are also scrolled through the display
  \item During the tuning procedure the current status is shown        
\end{itemize}

\section{Panel control}

(this needs to be done)

\section{Keyboard and performance section}

(this needs to be done)

\section{Key pad and parameter/settings menus}

With the firmware upgrade the Prophet 600 provides additional patch parameters and instrument settings. These are "hidden" in two new menus. The two menus are accessed using the \textbf{Number Pad} in the following way. 

\textbf{Additional patch parameters}

Note that the Prophet operates in two basic modes (live mode and preset mode, see section \ref{uimode} and the access to the additional patch parameters is different as commented in the text.

The additional patch parameters are stored with a patch. In live mode these parameters are accessed by pressing the buttons 0-9 (or pressing them twice or in some cases three times). The respective parameter name is scrolled through the display. The display then shows the value which is either numeric (shown) or a choice setting (scrolled through the display). The value change is done using the \textbf{Speed Dial} on the \textbf{Key Pad}. Important: in preset mode the display shows the patch (preset) number and by default the number pad is reserved for selecting the patch in this case. In order to access the additional patch parameters in preset mode the button \textbf{To Tape} must be activated (its LED will be lit).
       
The list of additional patch parameters can be found in section \ref{patchref}. Each parameter is explained in the respective functional context in the next chapter.

\textbf{Shift mode: miscellaneous settings and transposition}

The miscellaneous settings are patch independent but they are also stored, e.g. will be loaded the after switching the instrument off and on. The settings are accessed via the \textbf{Shift} mode. In this mode, different keys and buttons have a different function. To use this mode either hold down the button \textbf{From Tape} while selecting one of the additional functions or double click \textbf{From Tape} to enter \textbf{Shift Lock} mode (permanent shift) indicated by a blinking LED. To exit the mode push the button again. In shift mode the numbers 0...9 on the number pad access the miscellaneous settings (see below). The keys on the keyboard are used to set the keyboard transposition. 

Note that in contrast to the additional patch parameter pages, the miscellaneous settings associated with the numbers 0...9 are edited by pressing number key repeatedly to cycle through values or to initiate actions. The speed dial is not used. The first time a number is pressed (while holding From Tape) the display scrolls the current value or mode. Pressing the same number key again cycles through the values of that setting or confirms a certain action or mode change. The list of miscellaneuous settings can be found in section \ref{settingsref}. Each setting is explained in the respective functional context in the next chapter or in the appropriate section in the setup chapter.

\section{The work flow concept}

With the clearly laid out panel controls the the Prophet 600 is both, a hands-on performance instrument and sound sculpting synthesizer. The upgrade preserves that basic character. The controls on the panel have largely have the same function as in original except for some minor changes. However, the new functionalities of the upgrade required additional parameters and these had to be placed in the new menus. This has disadvantages. Diving into menus interrupts the work flow. It is abstract and unsatisfying. It is only acceptable for parameters which are rarely used and certainly not suited for live tweaking. Another improtant aspect is the perservation of the user interface concept. In particular, for each function there should be a unique and self-contained way to change it, instead of have a mix of say, dials and menues. This is one reason why he upgrade also contains "only" one LFO and two envelopes even though one might have added more with the new processor: this corresponds to the layout of the Prophet 600 and anything more would have meant at least a partial dive into menues. The upgraded Prophet 60 fulfills this ideal to a large degree but not fully. The LFO shape, for instance, is a mix of flipping a swithc on the control panel and choosing options in the additional patch parameter menu. However, this is one of a few notable exceptions.

A useful work flow of the upgraded Prophet 600 for creating or editing patches could be as follows:

\begin{enumerate}
  \item Pre-production: Set basic overall patch parameters, such as envelope shapes and routing, LFO targets and range 
  \item Production: Use the panel controls to interactively design the sound 
  \item Post-production: add additional features, such as vibrato and bend modulation for a final touch
\end{enumerate} 

Ideally, menu diving should be largely restricted to steps 1 and 3 so that uninterrupted hands-on and interactive work with the synthesizer is possible in step 2. It may not always be possible to finalize step 1 upfront and the user may occasionally have to go back to step 1.Nevertheless, the workflow should be more or less natural.

\section{Inputs and outputs}

\textbf{Power Switch}

When turning on the Prophet 600 with the upgrade installed the version is scrolled through the display. Then, the display will show and remain on the patch number selected when the instrument was last powered off.

Note: some earlier versions immediately went into tuning mode (see section \ref{tuning})  when powered up. From version 2.1 on the OS does not. This is to allow the instrument to warm up and stabilize, which typically takes approximately 10 to 15 minutes. Warm up time can depend on changes in air temperature, humidity and transportation of the instrument, as well as the condition of hardware such as the power supply. Therefore, after powering up, wait at least 5 to 20 minutes, and then, as needed, the TUNE button can be pressed to ensure optimal performance. If practical, allowing the instrument to warm up will result in less need to run the tune procedure and overall longer tuning stability.

\textbf{Line voltage selector}

The selector next to the main plug provides the option to switch between 110V and 220V mains voltage, whichever is appropriate for the line voltage where the instrument is operated. A flat head screw driver can be used to rotate the line voltage selector to the appropriate setting. Instrument should be powered down and unplugged when line voltage selector is adjusted.

\textbf{Fuse holder}

Fit the fuse according to the selected line voltage, e.g. 110V 1/2A, slo-blo, 220V 1/4A slo-blo.

\textbf{Audio out / headphone jack}

To drive a preamp or amp, a standard monophonic cable can be used. Or if there are two audio destinations (perhaps one to remain “dry” while the other is processed), a stereo cable may be more convenient.

The Prophet 600 has a monophonic output signal, but the jack is wired so that both sides of standard stereo headphones can be driven. The headphones have a minimum impedance of 1200 Ohms per element (600 Ohms in parallel).

\textbf{Filter CV In}

This jack accepts a 0-10 Vdc control voltage (CV) which raises programmed filter frequency settings. This enables remote and spontaneous increase (but not decrease) of brightness. This CV is usually provided by an accessory voltage pedal.

\textbf{Foot switch}

The foot switch jack allows for the use of a sustain pedal (only supported from GliGli 2.1 onwards). The pedal should be contact open at rest (for example a Fatar PS100 will function correctly).

\textbf{Cassette In/Out (disabled with GliGli mod)}

In the upgrade version there is no way to receive or send patch data via Cassette In/Out. Cassette In can be used for sending a pulsed clock to the Prophet 600 to sync the sequencer and/or arpeggiator.

\textbf{MIDI In/Out}

(this needs to be done)

\chapter{Functionality}

\section{Operating modes: preset and live mode}\label{uimode}

As with the original Prophet 600 firmware, there are two modes of operation, a \textbf{Preset} mode and a \textbf{Live} mode. The user can change between the two by pressing the \textbf{Preset} button. 

The idea of the live mode is that the patch that is playing corresponds to exactly the parameters as currently set by all dials, switches and additional patch parameters. The synth produces the sound as set and seen. In the preset mode, in contrast, the patch parameter as stored in the patch memory are applied. The synth produces the the sound of the preset. Still, once a dial or a switch on the panel or an additional parameter is changed, this change will take effect (note that this may lead to discontinuous changes if for example the current position of a knob is far from the  value stored in the patch). In this case  "mixed state" applies, partly patch (untouched controls) and partly panel controls (touched controls). The user can then store the current sound in the same patch or a different patch. The patch will be stored as heard. 


\section{Loading and storing presets}\label{loadstorepatches}

Patches are stored in \storagemode. To enter this mode from \presetmode or \livemode press the \record button, whose LED starts blinking. The Prophet 600 is then waiting for the user to enter the (two-digit) target patch number on the \termnumberpad. Once the second digit is entered, the patch is stored automatically and the LED of the \record button deactivates. 

Loading patches is done by entering \presetpatch and then selecting the patch on the \termnumberpad. Once selected there is a short confirmation on the display.

While the Prophet 600 waits for patch number entry (e.g. \storagemode or \presetpatch) other functions of the \termnumberpad and display (e.g. in \presetpanel or \livemode) are temporarily suppressed. If you find you have accidentally started typing a digit in these modes while expecting to select a patch parameter you can switch / cancel the mode after the first digit, e.g. pressing \totape in \presetpatch or pressing \record again in \storagemode.

For exporting and importing patches via MIDI SysEx see section \ref{patchmgmt}.


\section{Oscillator and mix}\label{osc}

Each of the six voices of the Prophet 600 has two voltage controlled oscillators (VCOs). VCO A and VCO B have dedicated sub-panels. The oscillators are basically identical. They offer three waveforms: \textit{sawtooth}, \textit{triangle} and \textit{pulse}. The pulse width of the pulse shape can be adjusted using the corresponding \pulsewidth dial. All three shapes can be activated simultaneously. The figures below are for \textit{GliGli} panel layout and \textit{SCI} panel layout, which differ in the way the oscillator volumes are controlled.

\textbf{Oscillators in GliGli panel layout}

\begin{center}
\scalebox{0.4}{
  \begin{tikzpicture}[scale=0.8]
    \addpar{-18cm,13cm}{\oscpitchmode \\ This menu parameter oscillator sets the units of the two oscillator frequency dials, e.g. \textit{free}, \textit{semi}, \textit{octave}, or mixed (different units for A and B.) };
    \oscapanelsci{0cm,9.5cm}
    \oscbpanelsci{0cm,0cm}
    \mixerpanel{31.9cm,0cm}
    \draw[line width = 2pt, dashed] (-4cm,15cm) -- (1cm,15cm);
    \draw[line width = 2pt, dashed] (-4cm,15cm) -- (1cm,6cm);
  \end{tikzpicture}
}
\end{center}

\textbf{Oscillators in SCI panel layout}

\begin{center}
\scalebox{0.4}{
  \begin{tikzpicture}[scale=0.8]
    \addpar{-18cm,13cm}{\oscpitchmode \\ This menu parameter sets the units of the two oscillator frequency dials, e.g. \textit{free}, \textit{semi}, \textit{octave}, or mixed (different units for A and B.) };
    \draw[line width = 2pt, dashed] (-4cm,15cm) -- (1cm,15cm);
    \draw[line width = 2pt, dashed] (-4cm,15cm) -- (1cm,6cm);
    \oscapanelsci{0cm,9.5cm}
    \oscbpanelsci{0cm,0cm}
    
    \prophetpotbiplar{35cm, 14.5cm}{MIX}{75}
    \node[rectangle, font=\fontsize{14}{12}\selectfont, anchor = center] at (33cm,12.5cm) {\panelfont{OSC A}};
    \node[rectangle, font=\fontsize{14}{12}\selectfont, anchor = center] at (37cm,12.5cm) {\panelfont{OSC B}};

    \addpar{42cm,13cm}{\drive \\ This menu parameter modifies the drive factor of the oscillator mix dial (applies to SCI panel layout only)};
    \draw[line width = 2pt, dashed] (41cm,15cm) -- (38cm,15cm);
  
  \end{tikzpicture}
}
\end{center}

\textbf{Oscillator controls}

Both oscillators have a \oscfreq dial which determines the frequency range of the oscillators when played from the keyboard or via MIDI.  You can choose the granularity of the frequency dials: 
\begin{itemize}
  \item \textit{free}: the base frequency can be chosen freely. The dial value selection is continuous and the \display shows approximate values between 0 and 99 in the normal way
  \item \textit{semitone}: the base frequency can be selected in semitones. The dial value selection is discrete and the \display shows values from 0 to 63 (semitones).   
  \item \textit{octave}: the base frequency can be selected in octaves. The dial value selection is discrete and the \display shows values c0, c1, c2, c3, c4 or c5. 
\end{itemize}  
In addition there are mixed settings, e.g. \textit{octave-semitone}, \textit{semitone-octave}, \textit{octave-free} and \textit{free-octave} providing different granularity for oscillators A and B. The parameter to select the granularity is the additional patch parameter \oscpitchmode. 

The additional \freqfine dial for oscillator B is bipolar and provides frequency fine tuning with a range of about $\pm 1$ semitone. The main purpose of separate frequency settings per oscillator is to be able to detune them against one another. This can be used to create thick and lively patches (when the pitch difference is small), to create complex timbres with harmonics (when pitch difference is large, typically octaves) or even create intervals (for example fifths). The frequency setting is also an important aspect of poly-mod, see section \ref{polymod}. To be able to tune the relative frequency very accurately, \freqfine has a very sensitive (e.g. slow changing) mid region between -1 and 1.

The oscillator A has an additional \oscsync switch. The hardware of the Prophet 600 offers hard sync of oscillator A to oscillator B, e.g. when \oscsync is activated, oscillator A is reset to the cycle start when oscillator B completes a cycle. 

Finally, the outputs of both oscillators are mixed into one signal before entering the filter and output stages. The volume of oscillators A and B are either set by the \vola and \volb dials in GliGli panel layout or the \mixer and \drive parameters in SCI panel layout as described in the following. For details on switching the panel layout, see section \ref{panelswitch}.

\textbf{Oscillator mix in GliGli panel layout}

The panel schematic is shown above. In GliGli panel layout the volume of each oscillator has a dedicate control in the mixer sub-panel. The normalised output of either oscillator is about half way of the respective dial. Beyond this, the amplifiers into the filter can be over-driven. This is a new feature of the upgraded firmware from version 2.0 on.

\textbf{Oscillator mix in SCI panel layout}

In the SCI panel layout the two dials are designated \mixer (becoming a bipolar dial) and \glidepot and these control the oscillator mix and the glide time, respectively. (For glide amount see section \ref{glide}.) In order to cover the entire range of oscillator volumes which the GliGli controls provide, there is an additional menu parameter \drive which becomes available available when the SCI panel layout is selected. The original Sequential Circuits mix function corresponds to a half-way value for \drive at which oscillators A and B are mixed linearly. At \mixer =-50 oscillator A is at full volume and oscillator B is silent. At \mixer =+50 oscillator B is at full volume and oscillator A is silent.  For smaller and larger values of \drive the oscillator mix function is modified as shown below.


\scalebox{0.4}{
  \begin{tikzpicture}[scale=0.8]
    \addpar{-2cm,14cm}{\textbf{Oscillator mix function} \\ In SCI panel layout the volumes of oscillator A and B are set by using the \mixer control and the menu parameter \drive.};

    \coordinate(orig)at(20cm,0cm);
      
      % A-axis
    \draw[-to, line width = 2pt] (orig) -- ($(orig)+(12cm,0cm)$); 
    \node[rectangle, font=\fontsize{17}{12}\selectfont, anchor = west] at ($(orig)+(12.2cm,0cm)$) {Osc B}; 
    \node[rectangle, font=\fontsize{15}{12}\selectfont, anchor = north] at ($(orig)+(10cm,-0.2cm)$) {Full}; 
    \node[rectangle, font=\fontsize{15}{12}\selectfont, anchor = north] at ($(orig)+(5cm,-0.2cm)$) {Half}; 
    \draw[-to, line width = 2pt] (orig) -- ($(orig)+(0cm,12cm)$); 
    \node[rectangle, font=\fontsize{17}{12}\selectfont, anchor = south] at ($(orig)+(0cm,12.2cm)$) {Osc A}; 
    \node[rectangle, font=\fontsize{15}{12}\selectfont, anchor = east] at ($(orig)+(-0.2cm,10cm)$) {Full}; 
    \node[rectangle, font=\fontsize{15}{12}\selectfont, anchor = east] at ($(orig)+(-0.2cm,5cm)$) {Half}; 

    \prophetpotbiplar{37cm,0cm}{MIXER}{310};
    \prophetpotbiplar{20cm,17cm}{MIXER}{240};
    \prophetpotbiplar{32.5cm,12.5cm}{MIXER}{90};

    \draw[to-to, line width = 2pt] ($(orig)+(10cm,0.2cm)$) -- ($(orig)+(0.2cm,10cm)$); 
    \node[rectangle, font=\fontsize{14}{12}\selectfont, rotate=-45, anchor = south] at ($(orig)+(4.2cm,6.2cm)$) {drive=50}; 

    \draw[to-to, line width = 2pt] ($(orig)+(5cm,0.2cm)$) -- ($(orig)+(0.2cm,5cm)$); 
    \node[rectangle, font=\fontsize{14}{12}\selectfont, rotate=-45, anchor = south] at ($(orig)+(2.4cm,3cm)$) {drive=25}; 

    \draw[to-to, line width = 2pt] ($(orig)+(2cm,0.2cm)$) -- ($(orig)+(0.2cm,2cm)$); 

    \draw[-to, line width = 2pt] ($(orig)+(7.6cm,7.6cm)$) -- ($(orig)+(0.6cm,9.9cm)$); 
    \draw[-to, line width = 2pt] ($(orig)+(7.6cm,7.6cm)$) -- ($(orig)+(9.9cm,0.6cm)$); 
    \node[rectangle, font=\fontsize{14}{12}\selectfont, rotate=-18.4, anchor = east] at ($(orig)+(7.7cm,8.04cm)$) {drive=75}; 

    \draw[-to, line width = 2pt] ($(orig)+(9cm,9cm)$) -- ($(orig)+(0.2cm,10.2cm)$); 
    \draw[-to, line width = 2pt] ($(orig)+(9cm,9cm)$) -- ($(orig)+(10.2cm,0.2cm)$); 
    \node[rectangle, font=\fontsize{14}{12}\selectfont, rotate=-6.39, anchor = east] at ($(orig)+(8.9cm,9.56cm)$) {drive=90}; 

  \end{tikzpicture}
}


\section{Volume and global tune}

To the right of the panel there two global parameter dials, the \textbf{Master Tune} and \textbf{Volume} dials. The master tuning has a range or approximately $\pm$semitone. It's pupose is to be able to tune the synth to other instruments without touching the internal, relative frequencies of the oscillators and thereby potentially changing the patch.

\section{Filter}\label{filter}

The mix of the two VCOs is routed through a resonant 24db low pass filter (built on a Curtis CEM3372 filter). The control for the filter consists of a \filtercutoff dial regulating the cut-off frequency, and a \filterres dial for the resonance. The \filterenv dial regulates how strongly the cut-off frequency is modulated by the filter envelope. This dial is bipolar, e.g. it supports positive and negatives values (effectively an inverse envelope setting). The display will show values between -50 and 50 depending on position. This is major enhancement compared to the original Prophet-600 firmware which  featured only positive values.

\begin{center}
\scalebox{0.4}{
  \begin{tikzpicture}[scale=0.8]
    \addpar{-18cm,9.8cm}{\textbf{Filter Limit} \\ Using the miscellaneous setting options the cut-off frequency can be limited to 22kHz, see section \ref{limitsett} for details.};
    \draw[line width = 2pt, dashed] (-4cm,9.3cm) -- (1cm,9.3cm);

    \addpar{-18cm,2cm}{\textbf{\filenv} \\ This modifies the overall shape and speed range of the envelope, e.g. it influences the attack, decay and release phases, see \ref{envelopes} for details};
    \draw[line width = 2pt, dashed] (-4cm,3.5cm) -- (1cm,3.5cm);

    \filterpanel{0,0}
  \end{tikzpicture}
}
\end{center}

The \keyboardtrack switch has the setting \textit{Off}, \textit{1/2} and \textit{Full}. It lifts the cut-off frequency depending on keyboard note, e.g. the cut-off frequency is higher for higher notes. The effect applies to the cut-off frequency resulting from the fixed \filtercutoff setting plus any envelope and LFO modulation, e.g. it applies to the sum.  

\textbf{Equal tempered tuning of the filter cut-off frequency}

In the Prophet-600 the filter cut-off frequency is tuned in the same way as the oscillator frequencies. Furthermore \keyboardtrack is implemented in a way to allow equal tempered play using the filter:

\begin{itemize}
  \item When \keyboardtrack is set to \textit{half} then the filter cut-off is increased by a semitone every two semitones on the keyboard 
  \item When \keyboardtrack is set to \textit{full} then the filter cut-off is increased by a semitone with every semitone on the keyboard, e.g. if follows the scale 
\end{itemize}

This feature essentially provides a third oscillator per voice. You can either play harmonically using the filter at self resonance even without oscillator A or B active (you should set the \filterenv dial to zero, at least in first experiments). The \filtercutoff acts as a tuning dial in this case. Or you can use the self resonant filter in combination with oscillators A and B if you "tune" the filter frequency appropriately. You can even use cross-modulation using the poly-mod function to modulate the filter cut-off using oscillator B (see section \ref{polymod}). However, there is no way to control the volume of the self-resonance relative to the volume of oscillators A and B.   


\section{Envelopes and routing}\label{envelopes}

As in the original version the upgraded Prophet 600 offers two ADSR envelopes. The first is the filter envelope, which is associated with and routed to the filter cut-off frequency. Its parameters are located in the filter section of the panel. The second envelope is assignable and the parameters have a dedicated Assignable Envelope section on the panel below the filter section \footnote{In the original Prophet 600 as well as in the uprage up to version 2.1 RC3 the assignable enevelope is in fact a dedicated amplitude envelope. Therefore, in the original and in the upgrade compatible panel overlay the panel section of this envelope is marked "Amplifier". Also with the change from version 2.1 to the version 2.2 the additional patch parameter Amplitude Shape has been consequently renamed Assignable Shape.}.

Both envelopes work identically in that they have a classic ADSR setup with an \textbf{Attack} dial, a \textbf{Decay} dial, a \textbf{Sustain} dial and a \textbf{Release} dial. Compared to the original, the upgraded Prophet 600 provides not only faster and smoother envelopes but also offers new two different envelope speed regimes, \textit{slow} and \textit{fast}, as well as two different envelope shapes, \textit{linear} and \textit{exponential}.

The shape and the speed can be set independently for the filter and the assignable envelope using two additional patch parameters. The parameter \textbf{(5) Assignable Shape} changes the assignable envelope and the parameter \textbf{(55) Filter Shape} changes the filter envelope. Each of the two parameters can take the following values:

\begin{itemize}
  \setlength\itemsep{0cm}
  \item \textit{slow -linear}
  \item \textit{slow - exponential}
  \item \textit{fast - linear}
  \item \textit{fast - exponential}
\end{itemize}
 
In general the envelopes have three potential targets: filter cut-off frequency, amplitude and poly-mod (for this see section \ref{polymod}). The filter envelope is always routed to the filter cut-off frequency and the modulation strength is controlled by the \textbf{Filter Envelope} dial. For the  amplitude and poly-mod targets different routing options exist as described in detail in the section on poly-mod \ref{polymod}. The additional patch parameter \textbf{(555) Envelope Routing} envelope routing is used to select this. The main benefit from the routing options is to make poly-mod independent from the filter dynamics. The default routing is that the poly-mod is target by the filter envelope and the amplitude is modulated by the assignable envelope.

The modulation strength for the amplitude is fixed in the Prophet 600 (e.g. there is no offset) and positive. The modulation strength for the poly-mod is set by the \textbf{Poly-Mod Envelope} dial as explained in section \ref{polymod}. 


\section{Polyphonic, unison and chord mode, voice assignment priority}\label{poly-unison-voice}

The upgraded Prophet 600 can be played in the following voice modes
\begin{itemize}
  \setlength\itemsep{0cm}
  \item \textit{polyphonic}: This mode is set by setting the \unison on the panel to \textit{off}. In this mode any new note is assigned to one of the 6 voices. 
  \item \textit{unison}: This mode is activated by switching the panel \unison switch from \textit{off} to \textit{on} without holding any keyboard key pressed. In this mode all 6 voices play the same note. Apply the additional patch parameter \detune to create thickened sounds in unison mode, see section \ref{detune}.
  \item \textit{chord}: This mode is activated by switching the panel \unison switch from \textit{off} to \textit{on} while holding down one or more keys (the "chord"). In this mode the synthesizer transposes that chord over the whole keyboard. Note that independently of how many notes the chord contains (and therefore how many voices would be left in principle), only one chord can be played at a time, similar to unison mode.
\end{itemize} 

The setting of \unison is a patch parameter e.g. is stored with the patch.

Note that in unison/chord voice modes the foot switch does not hold/sustain notes as in polyphonic mode but instead latches patterns of notes / chords on the fly without leaving and reactivating the unison mode. To do so, hold down the chord you want to latch in unison/chord mode and press the foot switch. In order to return to normal unison mode press the foot switch with no key held down. Using the foot switch with \unison on re-triggers the currently held notes. 

The assignment of voices to new played notes can be customized using the additional patch parameter \prio. New notes will be assigned to voices according to the following priority rules:

\begin{itemize}
  \setlength\itemsep{0cm}
  \item \textit{last}: New notes always play and if 6 voices are already playing then the oldest played notes is "stolen".
  \item \textit{low}: A new played note only plays if it is lower than the highest note already playing. The highest note is "stolen". With \unison on, legato is active. This is the preferred setting when bass accompaniment is important and the sudden stop of bass notes would create unwanted disruptions. 
  \item \textit{high}: A new played note only plays if it is higher than the lowest note already playing. The lowest note is "stolen". With \unison on, legato is active. This is the preferred setting for a synthesizer lead situation where the sudden stop of a high note would create unwanted disruptions. 
\end{itemize}

The assigner priority of voice mode is a patch parameter e.g. is stored with the patch.

\textbf{Ideas for solo voices}

The assigner priority can be used effectively in combination with the voice modes to customize solo voices. For example combine the priority \textit{high} (or \textit{low}) with ...
\begin{itemize}
  \item ...\textit{unison} mode for a massive legato solo voice
  \item ...\textit{chord} mode with a one key "chord" (hold down only C3) to achieve a more delicate legato mono solo voice
  \item ...\textit{chord} mode with octaves (hold down C3 plus C4 plus more if desired) for a more assertive mono solo voice   
\end{itemize}

In the default setup the voice allocation follows a "first voice found" logic. This means that when a new key is played the first free voice (first in the sense of the internal numbering of voices) will be allocated. In practise this means that when a note is released and played again the same voice is used. Since voices can have different characteristics in analogue synthesizer such as the Prophet 600, such a voice allocation logic favours more stable and consistent note sequences. However, the variations in voices can also be deliberately used to create more lively and animated sequences. This becomes very pronounced when vintage spread is activated (see section \ref{spreadsett}). In order to give users the choice, the advanced menu parameter \assign has been introduced which offers the choices \textit{first} and \textit{cycle}. The second option activates a "round robin" logic which always selects a new (free) voice for allocation, e.g. distributing the assignments evenly among the 6 voices.   


\section{LFO}\label{lfo}

As the original model, the upgraded Prophet 600 offers a single LFO (apart from vibrato, described in scetion \ref{vib}). However, the functionality is substantially enhanced. The frequency range of the LFO is much wider, from 1/20 to 60 Hz. There are additional wave forms (a total of 6). The upgraded Prophet 600 offers more modulation target options and, finally, an LFO delay can be set. These enhancements require 4 additional patch parameters as described in the following. However, apart from these parameters the most important LFO controls can be found and be changed hands-on in the dedicated LFO panel section. 


\begin{center}
\scalebox{0.4}{
  \begin{tikzpicture}[scale=0.8]
   \lfosubpanel{0,0}
  \end{tikzpicture}
}
\end{center}

\textbf{LFO speed}

The LFO speed / frequency is set in the panel using the \textbf{LFO Frequency} dial. In order to provide fine tuning on a broader frequency there are two basic frequency ranges, \textit{slow} and \textit{fast}. This is toggled using an additional patch parameter \textbf{(22) LFO Range}. 

\textbf{LFO modulation amount and targets}

There are three basic modulation targets which can be active simultaneously. Each one is activated using the switches \textbf{VCO Frequency}, \textbf{Pulse Width} and \textbf{Filter Frequency} in the LFO panel section. The modulation amount is set using the \textbf{Initial Amount} dial and this amount applies to all activated modulation targets. 

In the upgrade Prophet 600 the target VCO Frequency can be cutomized in the following way. The user can choose which of the two oscilliators per voice is modulation. The additional patch parameter \textbf{(11) LFO Target} offer the settings: \textit{VCO A \& B} (both osciallators are modulated, this is the standard of the original Prophet 600), \textit{VCO A} (only VCO A ist modulated), \textit{VCO B} (only VCO B ist modulated),  and \textit{VCO A \& B \& VCA} (both osciallators and amplitude is modulated). The targets pulse width and filter frequency are fixed, e.g. they can only be switch on or off and there are no further parameters for customization.

\textbf{LFO shapes}

The upgraded Prophet 600 supports not only the waveform \textit{Pulse} and \textit{Triangle} of the original model but four additional shapes, \textit{Sin}, \textit{Sawtooth}, \textit{Random} and \textit{Noise}. On the panel the is a simple up/down \textbf{Shape} switch. To accommodate 6 waveforms an additional patch parameter \textbf{(1) LFO Shape} is used. Selecting the desired waveform therefore requires selecting the respective shape pair and the using the \textbf{Shape} switch to change between the two shapes. 

\begin{tabular}{ |l|l|l|} 
 \hline
  \textbf{(1) LFO Shape} set to: & \textbf{LFO} switch \textit{up} & \textbf{LFO} switch \textit{down} \\
 \hline
 \textit{tri-square} & LFO shape is triangle & LFO shape is square\\
 \hline
 \textit{sin-rand} & LFO shape is sin & LFO shape is random \\
 \hline
 \textit{saw-noise} & LFO shape is sawtooth & LFO shape is noise \\
 \hline
\end{tabular}

\textbf{LFO delay and modulation of LFO}

As part of the patch parameters, the strength LFO effect itself can be modulated by either the modulation wheel or a delay function. The additional patch parameter \textbf{(44) Modulation Wheel Target} can be set to either \textit{vibrato} or \textit{LFO} (or \textit{off}. When set to \textit{LFO} the modulation wheel controls the LFO amount, e.g. the LFO starts at initial amount as set using the \textbf{Initial Amount} dial on the panel and turning the modulation wheel up adds LFO strength. The wheel effect itself is set using using the additional patch parameter \textbf{(33) Modulation Wheel Range}. If the modulation wheel target is set to \textit{vibrato} then, automatically, the the modulation delay is applied to the LFO, e.g. the LFO onset is delayed by the amount set by the additional patch parameter \textbf{(4) Modulation Delay}. For a details on moudlation wheel related parameters see section \ref{modwheel}.

\textbf{LFO Sync}

As a new feature in version 2.2 the LFO can be sync'ed to the arpeggiator / seequencer. The additional patch parameter \textbf{(111) LFO Sync} provides choices \textit{off} (free running LFO) and \textit{1, 2, 3, 4, 5, 6, 8} to have the LFO reset after 1, 2, 3, 4, 5, 6 or 8 beats. 


\section{Poly-Mod}\label{polymod}

The functions associated with the poly-mod section are very powerful. They provide tools for non-linear modulation: cross-modulation of the two voice oscillators and modulation of the filter cut-off frequency by VCO B. 

A familiar and common synthesizer feature is that oscillators and/or filters (frequency, filter cut-off frequency, pulse width, etc.) are modulated by a \textit{low frequency oscillator (LFO)}. An LFO does not operate in the audible frequency range (it is "low frequency") and it has a fixed frequency which is independent from the note played. Poly-mod is different in that the modulation itself is in the audible range and the resulting non-linearity changes the frequency spectrum drastically. The resulting effect is comparable to FM synthesis known from digital synthesizers, although in the Prophet-600 it is more raw and fully analog. Like the oscillator sync function, cross-modulation is hardware implemented in the Prophet-600.

The poly-mod control can be found in the poly-mod sub-panel. 

\begin{center}
\scalebox{0.4}{
  \begin{tikzpicture}[scale=0.8]

    \addpar{-18cm,3.5cm}{\envrouting \\ The routing options determine whether poly-mod is driven by the filter or the 2nd envelope};
    \draw[line width = 2pt, dashed] (-4cm,3.5cm) -- (1cm,3.5cm);


   \polymodesubpanel{0,0}
  \end{tikzpicture}
}
\end{center}

There are two modulation sources, \polymodenv and \polymodosc and two destinations, the pitch of oscillator A  and the filter cut-off frequency. Each source is activated using the respective panel switch, \polymodfreq and \polymodfilter. The \polymodosc  determines how strongly the oscillator B modulates the poly-mod targets (destinations). It is instructive to consider the effect of oscillator B modulation on the two targets first and then explore the effect of envelope modulation.

\textbf{Cross-modulation}

\begin{itemize}
  \item When you set the \polymodfreq switch to \textit{on}, oscillator B modulates the frequency of oscillator A (frequency modulation, FM). Depending on the modulation strength the result is a wide and rich frequency spectrum of oscillator A which contains among others the sum and difference of multiples of the two oscillators frequencies. The spectrum strongly depends on the frequency difference and may be totally non-harmonic but also harmonic for carefully adjusted frequency settings. Typical spectra are often described as bell sounding. Note that in contrast to FM synthesizers which use highly accurate sin waves, the oscillator output of the Prophet-600 does not have the same basic accuracy and the available waveforms already contain higher harmonics (triangle, sawtooth, etc.). Therefore the FM result in the Prophet-600 is typically more drastic and dirtier and requires classic subtractive methods for shaping the sound, i.e. the application of the low pass filter. 
  \item When you set the \polymodfilter switch to \textit{on}, oscillator B modulates the filter cut-off frequency of the VCF. Like frequency modulation this is a non-linear modulation. This effect can also produce a wide frequency spectrum which may or may not be harmonic, depending on the relative frequency of oscillators A and B. The effect is more pronounced at low filter cut-off frequencies. The resonance also has a strong influence on the effect.  
\end{itemize}

Note that poly-mod is very pitch sensitive so that patches which use strong poly-mod may need to be fine tuned using \freqfine. Also, poly-mod patches may sound off on a cold unit due to temperature dependent drift and generally it maybe necessary to retune the synth more often when using extreme poly-mod patches.

\textbf{Envelope pitch modulation}

The second poly-mod modulation source is the poly-mod envelope. When the \polymodfreq switch is \textit{on}, the envelope modulates the pitch of oscillator A at the strength set by the \polymodenv dial. The \polymodenv control is bipolar. It  provides negative (inverse modulation, lowering the pitch) and positive modulation (raising the pitch). The pitch shift is harmonic and in the sustain phase of the envelope the maximum obtainable pitch shift (i.e. maximum positive or negative value of \polymodenv and sustain at full value) corresponds to $\pm 30$ semitones.

Note that the poly-mod envelope only affects the frequency of oscillator A (and never the filter cut-off frequency) even if the second destination, the \polymodfilter switch, is also \textit{on}. The filter cut-off frequency is already modulated by the filter envelope. Depending on the menu parameter \envrouting the source envelope is either the filter envelope or the 2nd envelope, see below for details on the options.

Pitch modulation by envelope is not actually a cross-modulation function. However, the association with poly-mod is useful because of the strong effect pitch (difference) has on cross-modulation. What can it be used for? Envelope modulation is a form of time modulation. It can change the sound from attack through hold to release. By time modulating the pitch in combination with cross-modulation (and/or sync), the sound can be made to sweep through different spectra, ranging from apocalyptic noise through bell sounds to delicate transients. The way humans perceive natural sounds is strongly influenced by the frequency characteristics of the transients (typically short lived high frequency components) in the attack and decay phases of the sound. This is best modelled using an envelope. So, being able to modulate the composition of frequency spectrum by envelope is a versatile tool for advanced sound design. This was probably the original thought behind putting the \polymodenv control in the poly-mod section.

In order to further expand the sound design possibilities, the latest firmware upgrade offers an envelope routing option, which is inspired by a similar setup in the Roland JX8P synthesizer. Having different envelopes modulating the filter and the poly-mod opens up new possibilities in the sound and spectrum dynamics. The following section explains the routing options. 

The menu parameter \envrouting determines if the filter envelope or the 2nd envelope is used in poly-mod\footnote{Up to version 2.1 RC3 the filter envelope was "hard wired" to poly-mod.}.

\begin{itemize}
  \setlength\itemsep{0cm}
  \item \textit{standard}: In this "standard" routing the amplitude is modulated by the 2nd envelope and the poly-mod is modulated by the filter envelope. This is the setup of the original Prophet-600 and versions 2.0 and 2.1 RC3. It is also the default in \version.
  \item \textit{poly-amplitude}: In this routing both the amplitude and the poly-mod are modulated by the 2nd envelope. 
  \item \textit{poly}: In this routing the amplitude is modulated by the filter envelope and the poly-mod is modulated exclusively by the 2nd envelope. Therefore the poly-mod time modulation is completely free, at the expense of having filter modulation and amplitude tied together.
  \item \textit{gate}: In this routing the amplitude is not modulated at all, but behaves like a gate (similar to electronic organs). The poly-mod is modulated exclusively by the assignable envelope. Therefore the poly-mod modulation is completely free at the expense of a less dynamic sound evolution.
\end{itemize}  

\scalebox{0.4}{
  \begin{tikzpicture}[scale=0.8]
   \envroute{0,0}
  \end{tikzpicture}
}

\textbf{Oscillator sync and poly-mod}

As explained in section \ref{osc}, you can sync oscillator A to oscillator B using the \oscsync switch in the sub-panel for oscillator B. This can also be view as a variant of cross-modulation. The sync function can also be combined with frequency modulation as described above. 

Combining oscillator sync with poly-mod pitch modulation by envelope can be used to produce interesting short transients or slow sound evolutions. For example, when oscillator A is a pulse wave, then sync'ing this wave to oscillator B and time modulating the pitch of oscillator A has the effect of pulse width modulation by envelope. It creates a change in timbre which can be short to simulate plucking an brassy attacks or long to produce a sweep of the sound spectrum.

\textbf{Starting points for patch design with poly-mod and sync}

Using poly-mod in combination with oscillator sync, different modulations strengths for different oscillator pitches with different filter and modulation shapes - this offers a very wide field of experimentation. Here are two prototypical starting points:

\begin{itemize}
  \item Set the envelope routing to \textit{standard}, tune both oscillators to the same frequency. In the poly-mod section turn source VCO B to zero and deactivate the filter modulation. Then set VCO A to sync and switch on the poly-mod target frequency oscillator A. Experiment with the poly-mod envelope source amount, oscillator A frequency and filter envelope attack and decay. This produces an evolving sync sound which, for high envelope amounts, can be tuned to a very assertive synth solo sound.
  \item Set the envelope routing to \textit{standard}, tune both oscillators to the same frequency. In the poly-mod section turn the poly-mod envelope source amount to zero, deactivate the filter modulation  and switch on the poly-mod target frequency oscillator A. Switch off VCO A sync. Now experiment with the poly-mod envelope source amount and oscillator B fine tune. Either you enjoy the drastic spectra you will get for most of the settings or you carefully adjust the frequencies to find the sweet spots at which the spectrum clears. This is the typical way to produce very rich and brassy timbres but normally this also needs some filter cut-off to softly cut away the highest frequencies. For stronger poly-mod, oscillator B fine tune may not be sufficient and it may be necessary to also adjust the frequency of oscillator A. For this it is useful to set the granularity of the frequency dials to free/octave (i.e. oscillator A free and oscillator B octave). Notice that other controls also influence the pitch, such as the oscillator B wave selection and pulse width (see description of bleed bug below) so that you may find that you need to compensate for this effect with fine tuning.   
\end{itemize} 

\textbf{Technical excursion: pulse width bleed bug}

The firmware versions 2.0 and 2.1 RC3 contain a software error which incorrectly resets the pulse widths of oscillators A and B. This manifests itself under the following conditions: \polymodosc is turned up and the poly-mod target \polymodfreq is active. The wave shape pulse of oscillator B is not activated. In this case the incorrect reset of the pulse widths leads to a wrong frequency in oscillator A. The effect becomes stronger with higher settings of \polymodosc. In version \version this bug is corrected but in order to ensure backwards compatibility of older patches, the bug is automatically switched on when loading such patches to faithfully reproduce the behaviour of prior firmware versions.  The bug toggle is the \syncbug which can be set to \textit{off} and \textit{on}. The default value (default patch) is \textit{off}.

The Prophet-600 has a hardware error which leads to a bleed of pulse width voltage of oscillator B to the pitch of oscillator A when \polymodosc with destination \polymodfreq is active, e.g. when the pulse waveform of oscillator B is active, the pulse width of oscillator B changes the pitch of oscillator A. In practice \oscfreq and/or \freqfine need to be used to compensate for the drift.   


\section{Arpeggiator}\label{arp}

The instrument supports three arpeggiator modes.

\begin{enumerate}
  \item \textit{up/down} plays an up/down note sequence with double trigger on the lowest and highest note. For example, playing C3 C4 E4 G4 will play: C3 C4 E4 G4 G4  E4 C4 C3 C3 C4 and so on. Mathematically, it allows you to work out the time signature.  Three notes result in a 3/4 time, 4 notes in a 4/4 time and 5 notes in 5/4, etc. To start the arpeggitor in \textit{up/down}, press the key \textbf{Arpeg Up/Dn} on the key pad. Its LED will light solid when active.
  \item in the \textit{assign} mode the arpeggiator plays the sequence of notes in the order in which they are played. To start the the mode press the \textbf{Arpeg Assign} button on the key pad and play the notes in the order you want them sequenced. The key LED will light solid.
  \item in the \textit{random} mode the arpeggiator replays the played notes at random.  To enter this mode press the \textbf{Arpeg Assign} key twice until its LED blinks and then play the notes. 
\end{enumerate}

The foot switch input can be used to hold the arpeggiator. Notes send to the Prophet 600 via MIDI are not added to the arpeggiator. It is possible to sync the LFO to the arpeggiator using the additional patch parameter \textbf{(111) LFO Sync}, see section \ref{lfo}.

\textbf{Latch mode}

With all arpeggiator modes, press the \textbf{Record} on the key pad to enter the latch mode where played notes are held. The LED of the button will light solid. Playing additional notes in this mode will add them as additional notes to the existing sequence, up to a maximum of 128 notes. To clear the notes from the sequence, press the Record key to switch it off. The LED of the Record button will go off.

\textbf{Tempo}

The tempo of the sequencer is set by the \textbf{Speed} dial. To activate the tempo function the \textbf{(0) Seq/Arp Speed} parameter must be selected. Note that while in latch mode the button \textbf{To Tape} has to be activated (its LED is lit solid) to select additional patch parameters via the number pad and also to switch on the speed dial. The  areggiator can be synced to external sources. The miscellaneous setting \textbf{Clock Sync} (setting 8) pprovides the choices \textit{internal}, \textit{MIDI} or \textit{Tape In} (see also section \ref{sync}).

\textbf{Trouble shooting}

If the arpeggiator is activated and keys are held down or have been latched and still no sound plays, then it is worth checking the following potential reasons: Is the sync set to \textit{internal} and is the speed greater than zero (see settings in \ref{settingsref})? If the sync is not \textit{internal} does the instrument receive a clock, e.g. via MIDI on the correct channel or via tape in?


\section{Sequencer}\label{seq}

The live sequencer of the original Prophet-600 has been replaced by a step sequencer which supports polyphony. You can program two tracks which can play simultaneously. Sequence 1 is started by pressing the button \seqone, sequence 2 is started by pushing the button \seqtwo. The LED of the running sequencer is always on. Live sound control via panel controls, menu parameters and MIDI CC as well as playing via keyboard and/or MIDI is fully supported while the sequencer is running. To stop either sequence, press the respective sequencer button again. The LED of the stopped sequencer is off. 

By default the sequences are empty, so even with the started sequencer no sound will play. If either sequence is activated and no sound plays, then it is worth checking the following potential reasons: Has a note been added? Is the sync set to \textit{internal} and is the speed greater than zero? If the sync is not \textit{internal} does the instrument receive a clock, e.g. via MIDI on the correct channel or via tape in?

\textbf{Creating / editing a sequence}

The following is an example work flow for adding notes, rests and ties to a sequence that covers all editing functions. To enter the editing mode press the \record button and then the respective sequencer track button as shown below (example for editing sequencer track 1):

\begin{center}
  
\scalebox{0.4}{
  \begin{tikzpicture}[scale=0.8]
    \arpsqbuttons{9cm, 6cm}{}{}
    \draw [-to, line width=2pt](22cm,9cm) -- (28cm,9cm);
    \node[font=\fontsize{22}{22}\selectfont, align=left, outer sep=0.5mm, anchor = south west, text width=10cm] at (22cm,9.5cm) {Press \record};

    \arpsqbuttons{29cm, 6cm}{R}{R}
    \draw [-to, line width=2pt](42cm,9cm) -- (48cm,9cm);
    \node[font=\fontsize{22}{22}\selectfont, align=left, outer sep=0.5mm, anchor = south west, text width=10cm] at (42cm,9.5cm) {Press \seqone};

    \arpsqbuttons{49cm, 6cm}{R1}{1}
  \end{tikzpicture}
}
\end{center}

\begin{enumerate}
  \item Press the \record button. Its LED starts flashing.
  \item Press the sequencer button of the sequence you would like to edit (\seqone or \seqtwo). The LED of the respective sequencer button starts flashing, and the \record button LED becomes solid. The sequence is in \seqedit and awaits notes, ties and rest to be added to the sequence. Note: The sequencer is running in this mode, so as soon as the first note is entered this note is repeatedly replayed at the current sequencer/arpeggiator tempo. If notes have previously been added to the sequence, this sequence is played and any new notes are added. 
  \item Adjust the tempo. Note that you can set the \datadial to control the clock tempo before starting the sequencer editing mode. The parameter selection is preserved. Be careful when selecting the parameter in record mode. It requires that you press and hold down the \totape button. Otherwise pressing the \numberbut{0} has a different, sequencer specific function: it clears the sequence!  
  \item Play the notes to be added to the sequence. The display screen shows the number of steps already added, e.g. it is the current length of the recorded sequence including rests and ties.
  \item Pressing \numberbut{0} resets the sequence and deletes any notes. Note: there is no confirmation step. 
  \item In \seqedit the \numberbut{2} on the \termnumberpad inserts rests or ties. It adds a rest when no key is pressed, and a tie when a key is pressed.
  \item Added notes (or chords), ties and rests can be successively deleted/undone using the \numberbut{1} on the number pad. 
  \item In order to complete the sequence editing and to leave the \seqedit, press the respective sequencer button (\seqone or \seqtwo). The sequencer button LED becomes solid, and the \record button LED does off. The sequence is saved. 
  \item To start playing the sequence, press the respective button, (\seqone or \seqtwo). The respective LED is on when a sequence is running.
\end{enumerate}

\textbf{Transposing a  sequence}

To transpose a sequence use the keyboard in \shiftmode or \shiftlock. The C3 key is zero or standard tuning, and will not transpose. Playing C sharp 3 will transpose the sequence up a half step, playing D3 will transpose the sequence up a whole step, etc. See section \ref{transposition} for details.


\section{Vibrato}\label{vib}

The upgraded Prophet-600 features a dedicated vibrato. Its parameters are available in the additional patch parameter menu as follows.

\begin{itemize}
  \item The \vibspeed sets the vibrato frequency in a range of approximately 1/20 to 60 Hz
  \item The \vibamt sets the vibrato strength
  \item The \vibtgt sets the vibrato target. The options are \textit{VCO} (modulating the pitch of oscillators A  and B), \textit{VCA} (modulating the amplitude) as well as \textit{VCO A} and \textit{VCO B}(modulating only the pitch of one oscillator A  or  B, respectively)
\end{itemize}

The vibrato shape is a triangle wave. The vibrato effect can be controlled by either the modulation wheel or a modulation delay function. The additional patch parameter \modwheeltarget can be set to \textit{vibrato} or \textit{LFO}. When set to \textit{vibrato}, pushing the \modwheel up adds vibrato strength. The strength of the \modwheel is controlled by the additional patch parameter \modwheelrange. If you set the \modwheel target to \textit{LFO}, on the other hand, the modulation delay is automatically applied to the vibrato, e.g. the vibrato onset is delayed by the amount set by the additional patch parameter \moddelay. For setting the parameters of the modulation wheel see section \ref{modwheel}. For complementary functions for LFO see section \ref{lfo}.


\section{Pitch bender}\label{pitchbend}

The Prophet 600 comes with a (pitch) \textbf{Bender} (without automatic return to the middle position). With the upgrade it is possible to customize the range/action and the target of the bender as part of the patch confguration (see below). 

Note: External MIDI pitch bend is \underline{added} to the internal pitch bend\footnote{This is the behaviour of the original firmware. Since the bender does not return to zero a bend override by external MIDI leads to unwanted discontinuous pitch changes when the bender of the Prophet 600 is touched. The addition of both bends is therefore more consistent.}. Full internal and external (MIDI) bend covers twice the range set by the bender range. As a consequence, the external MIDI bend value cannot be reset/overridden using the built in bender, but the external value is reverted to zero during the bender calibration process (see below).

\begin{itemize}
  \item The bender can target the oscillator pitch, the filter cut-off frequency or the amplitude. This is set using the additional patch parameter \bentarget which provides the options \textit{off}, \textit{AB} (bend of the pitch of both oscillators), \textit{VCF} (modulating the cut-off frequency), \textit{volume} (modulating the amplitude) and \textit{b} (bending the pitch of oscillator B only).
  \item The range of the bender can be configured using the additional patch parameter \bendrange. It offers the choices \textit{2nd} (bend a full tone), \textit{3rd} (bend 3 semi-tones), \textit{5th} (bend 5 semi-tones) and \textit{octave} (bending an entire octave).
\end{itemize}

The upgraded Prophet 600 also offers the option to calibrate the bender middle position using one of the miscellaneous settings/function (see section \ref{settingsref}). This function is useful to remove jittery behaviour of the bender around the middle position. In order to calibrate the bender follow the steps below.

\begin{enumerate}
  \item Move the bender to the middle position.  
  \item Activate bender calibration by pressing \numberbut{3} in \shiftmode or \shiftlock (the display will show that the bender calibration mode is active). Pressing \numberbut{3} again applies and confirms the calibration and exits the calibration mode.
\end{enumerate}

As mentioned above, the current active external (e.g. MIDI) bend value is set to zero during the calibration process. To be precise, it is set to zero with the first press of \numberbut{3} in \shiftmode or \shiftlock.


\section{Modulation wheel and modulation delay}\label{modwheel}

The upgraded Prophet-600 has two methods for controlling the modulation effect of the LFO and the vibrato and these properties are  stored with the patch. The first one is the \modwheel. The target of the modulation wheel can be set using the menu patch parameter \modwheeltarget which provides the options \textit{LFO} and \textit{vibrato}. 

The second method is a modulation delay which delays the onset of the modulation effect. The delay is always applied to the modulation source which is \underline{not controlled} by the \modwheel, e.g. when the modulation wheel target is set to \textit{LFO} the delay applies to the vibrato and when it is set to \textit{vibrato} the delay is applied to the LFO.

The target setting as well as the modulation wheel strength and the delay time are set as follows. 
\begin{itemize}
  \item \modwheeltarget provides choices \textit{LFO} and \textit{vibrato}
  \item \moddelay is a continuous menu parameter (range 0...99) which controls the delay time
  \item \modwheelrange offers 4 options: \textit{touch}, \textit{soft} (this is the default), \textit{high} and \textit{full}
\end{itemize} 

The parameter \modwheelrange affects the \modwheel action for all variations except \textit{VCA}, e.g. if the target is \textit{vibrato} and the vibrato target is \textit{VCA} the wheel action is always effectively \textit{full}. The reason is that amplitude modulation is a weak effect and there is no practical value in making the \modwheel softer in this case.

Note: the \modwheel action has been modified in all settings compared to version 2.0 and 2.1 RC3. It is now exponential to provide a smooth onset for better playability.


\section{Glide}\label{glide}

In the original Prophet 600 the glide setting (the glide time) had a dedicated dial. In the upgraded Prophet 600 the dial has been re-designated to control the amplitude of OSC B. As a consequence the glide time has been moved to the additional patch parameter \textbf{(7) Glide}. 

\section{Detune}

The upgraded Prophet 600 provides a method to detune the 6 voices against one another. This effect can be used in particular to make unison sound broader and more lively. The additional patch parameter is \textbf{(77) Detune}

\chapter{Synthesizer setup}

\section{Tuning}\label{tuning}

The Prophet 600 is a VCO synthesizer, e.g. the frequency of the oscillators A and B are set by apllying a voltage to the control pin of the microchip. The relationship between absolute voltage and resulting frequency is not precise. It depends on many factors, such as smallest differences in the manufacturing process of the chip itself and the surrounding circuit components and operating temparture (the temparture of the components also changes during operation). Therefore the synthesizer need to be tuned regularly. Typically, tuning is best after the instrument (or rather: the chips and components) has been given time to warm up to a stable operating temperature, which could be 10-30 minutes into operation. It may be necessary to retune the synth during operation. The user should trust her or his ear on this.

In addition to a full tuning, each voice can be manually tuned. This per voice tuning is stored as part of the patch (in contrast to the thoverall tuning, which is stored as part of the settings which apply to all presets) an it is retained after full tuning. To enter the per voice tuning hold down From Tape and press Tune. The Tune button flashes. In this mode the last played note is tuned using the mod wheel. From the middle position the tune range is $\pm$1 semitone from equal tempered.

To leave the per voice tuning mode press Tune again.

\section{MIDI integration}\label{midiintegration.tex}

For supporting the integration of the Prophet-600 in a studio or live environment using MIDI, several MIDI functions are supported.

\begin{itemize}
  \item Note send and receive
  \item Pitch bend and modulation wheel actions
  \item Receive of all patch parameters using MIDI CC
  \item Local on and local off modes
  \item SysEx support for patch management (for this see section \ref{patchmgmt})
\end{itemize}

The MIDI send and receive channels can be set using dedicated miscellaneous settings as follows. In \shiftmode or \shiftlock:

\begin{enumerate}
  \setlength\itemsep{0cm}
  \item \numberbut{1} cycles through the MIDI receive channels, from \textit{OMNI} through channels \textit{Ch1} to \textit{Ch16}
  \item \numberbut{2}  cycles through the MIDI send channels, from \textit{Ch1} to \textit{Ch16}
\end{enumerate}

\textbf{Notes, performance and MIDI CC}

The Prophet-600 sends the following MIDI signals. 

\begin{enumerate}
  \setlength\itemsep{0cm}
  \item Key on, key off in all modes
  \item Pitch bender
  \item Hold pedal (in polyphonic mode, not in unison/chord mode or running arpeggiator or sequencer)
  \item Modulation wheel (technically a MIDI CC event)
  \item Program change (technically a MIDI CC event)
\end{enumerate}

All of the events can also be received and executed with the following exceptions/modifications.

\begin{enumerate}
  \setlength\itemsep{0cm}
  \item In arpeggiator mode incoming MIDI note on/off events are not entered into arpeggiator but play normally when in \textit{local on}. To enter notes to the arpeggiator using MIDI you need to switch to \textit{local off} mode. 
  \item In \shiftmode or \shiftlock incoming MIDI notes do no cause transpositions (but are ignored) in \textit{local on} mode. 
  \item External MIDI pitch bend is added to the pitch bend from the on-board bender 
  \item MIDI Hold Pedal events are also applied in unison/chord mode or running arpeggiator or sequencer where it has the effect of latching notes/chords
  \item MIDI program change messages are only applied in \presetmode and are ignored otherwise  
\end{enumerate}

The instrument does not send MIDI CC (apart from modulation wheel and program change). However, the Prophet-600 supports a comprehensive list of general and specific MIDI CC events. For reference see section \ref{midiimplementation}. Incoming MIDI CC events are only applied in \presetmode.

\textbf{Local off mode}

The Prophet-600 supports a \textit{local off} mode which is activated using \numberbut{0} with \shiftmode or \shiftlock. In \textit{local off} mode various internal events (such as notes) are not executed in the Prophet-600 directly but they are still sent as MIDI events. The typical setup would be to connect to a MIDI recorder and feed the recorded MIDI messages back to the instrument for replay. The \textit{local off} mode ensures that there is no doubling of events.

Note that in \textit{local on} mode external MIDI does not play into the arpeggiator but plays over the arpeggiator while the local keyboard plays into the arpeggiator. Since the arpeggiator and the playing the external MIDI notes share the 6 voices this can lead to interference of the two, such as missed arpeggiator notes, depending on what play using MIDI. When in \textit{local off} mode MIDI notes play into the arpeggiator (if active) and in this mode you can also use external MIDI notes for transposition of arpeggiator and sequencer like you would with the internal keyboard (i.e.. in \shiftmode or \shiftlock). However, transposition does never affect normal playing via MIDI in any mode. The reason is that the Prophet-600 sends MIDI notes with transposition already applied. For consistency reasons, incoming MIDI (which may in fact be a recording of MIDI notes from the Prophet-600) must therefore not be transposed again. If you need transposition from your external MIDI controller you would need to use whatever function the MIDI controlled provides to transpose the MIDI notes sent to the Prophet-600.

\begin{samepage}

\textbf{Overview of MIDI integration}

\nopagebreak
\renewcommand{\arraystretch}{1.3}

\begin{longtable}[l]{l|ccc|ccc}
     &
      \multicolumn{3}{c}{\textit{local on}} &
      \multicolumn{3}{c}{\textit{local off}} \\   \cline{2-7} 
    & Apply local & Send MIDI & Apply MIDI & Apply local & Send MIDI & Apply MIDI \\
    \hline
      \vline Keys on/off playing	& Yes & Yes	& Yes & No & Yes & Yes \\
      Keys into arpeggiator	& Yes & Yes	& No\footnote{MIDI notes will play normally / freely on top of the running arpeggiator. note that voice sharing between arpeggiator and playing MIDI notes can lead to unwanted interference such as missed arpeggiator notes etc.} & No & Yes & Yes \\
      Keys into sequencer	& Yes & Yes	& Yes & No & Yes & Yes \\
      Modulation Wheel	& Yes & Yes	& Yes & No & Yes & Yes \\
      Pitch Bend	& Yes & Yes	& Yes\footnote{Internal and external pitch bends are added, see \ref{pitchbend}} & No & Yes & Yes \\
      Program Change	& Yes & Yes	& in preset mode & Yes & Yes & in preset mode \\
      Pedal Hold (Latch)	& Yes & No & No & Yes & No & No \\
      Pedal Hold (Sustain)	& Yes & Yes & Yes & No & Yes & Yes \\
      Transpose & Yes & No & No & Yes & No & Yes\footnote{You can transpose the sequencer and the arpeggiator by MIDI, but this has no effect on normal MIDI playing.} \\
      Other MIDI CC	& Yes & No & in preset mode & Yes & No & in preset mode \\
\end{longtable}

\end{samepage}


\section{Voice deactivation}

The upgrade Prophet 600 provides a method for switching single voices of the Prophet 600 off. This can be useful to suppress a voice that is damaged or for troubleshooting, e.g. isolating voices for testing. The "voice masking" is done using the settings 4 and 5:
\begin{enumerate}
  \setlength\itemsep{0cm}
  \item Pressing key 4 while holding down From Tape repeatedly cycles through the 6 voice of the Prophet 600. This selection determines the functionality of the key 5 in the setting menu.
  \item Key 5 toggles between on and off of the voice selected using key 4. In the upgrade Prophet 600 the user can deactivate selected voices in this way.
\end{enumerate}

\section{Transposition}

The upgraded Prophet 600 provides a method for transposing the keyboard. Pressing key 7 while holding down From Tape activates a mode in which the user can set the keyboard transposition. The mode is indicated by a flashing From Tape LED. In this mode the last pressed key will define the transposition, the middle C3 being zero transposition. The value is also scrolled through the display for conformation. Pressing From Tape laves the transposition mode and the LED goes off.

Transposition also applies to arpeggiator and sequencer. The transposition procedure also works in the respective run, latch and record modes.

\section{Sequencer/arpeggiator sync}\label{sync}

The sequencer and arpeggiator of the upgraded Prophet 600 can be sync'ed to external clock sources. This is done using the miscellaneous setting of number key "8". Pressing key "8" repeatedly cycles through the sync source. The supported sources are \textit{internal} (default setting, the Prophet 600 generates its own clock.), \textit{MIDI} (fraction of the incoming MIDI clock on the received channel will be used as clock) and \textit{tape} (a pulse clock source can be plugged into the “Cassette In” jack and a fraction of it will be used as clock).

\section{Spread}\label{spreadsett}

(this needs to be done)

\section{VCF limit}\label{limitsett}

In the upgraded Prophet 600 it is possible to limit the cut-off frequency of the filter to 22kHz in order to avoid strange filter behaviour in the ultrasonic range.

In the miscellaneous settings menu this setting is combined with the spread setting (see section \ref{spreadsett} on spread).  

\section{Managing sound libraries via SysEx}\label{mididump}

The following instruction assume that you have your Prophet 600 connected via MIDI to a device where you can manage MIDI SysEx files, e.g. receive, store and send. Note that in the current firmware version the entrie patch bank of the Prophet is received or sent. Single patch transfer is not supported. 

\textbf{Dumping patches} 

On the settings menu the number pad "6" (use a second press to confirm) sends a SysEx dump of all presets to MIDI out on the MIDI send channel selected in settings.

\textbf{Loading patches} 

Loading a SysEx patch library will overwrite all existing stored patches. To load a complete SysEx Library:

\begin{enumerate}
  \item Make sure that the MIDI receive channel on the Prophet 600 and the MIDI send channel of the sending device are the same
  \item Transmit the data to the instrument   
\end{enumerate}

\section{SysEx firmware upgrade}\label{fwupgrade}

For a firmware upgrade in a Prophet 600 with an original Z80 processor please refer to the modification instructions \cite{modinstructions}. The following section is for owners of a modified Prophet 600 with Teensy board and a prior version of the GliGli upgrade installed. If modifying a Prophet 600 for the time one may directly flash the Teensy board with software version 2.2. It is not necessary to first install a lesser version and then upgrade.

Firmware upgrades are provided as MIDI SysEx files (*.syx) \footnote{They are also provided as *.hex files for USB transfer to the Teensy board.}. In order to perform the upgrade follow to step below.

\begin{enumerate}
  \setlength\itemsep{0cm}
  \item Switch off the instrument
  \item Switch on the instrument while holding the buttons \textbf{To Tape} and \textbf{From Tape} down. The display should show an "U" for upgrade on the left digit.
  \item Send the firmware SysEx file to the instrument. This process takes a while. The display shows a spinning segment. Once the process is completed, the display shows an "S" for success. If the upgrade ends with “E” for error on the display, the upgrade has failed and the upgrade procedure must be repeated until it succeeds. It is not recommended to start the instrument in normal mode after a failed upgrade.
  \item Switch the instrument off and on again. The scrolled welcome message should state the version "GliGli's P600 Upgrade Version \textit{<version>}"
\end{enumerate}

Users are recommended to follow the MIDI SysEx part of the modification instructions \cite{modinstructions}. In particular, the "Delay between buffers" and the "delay after F7" should be set to at least 250ms.

The settings and preset parameters are left unchanged by the upgrade, provided the previously installed version is one for which backward compatibility is supported. This applies to the following versions:

\begin{itemize}
  \setlength\itemsep{0cm}
  \item Version 2.0
  \item Version 2.1 Release Candidate 3 (not officially released)
  \item Version 2.2  
\end{itemize} 

Downgrading versions is possible but preservation of settings and preset parameters is not guaranteed in that case.

\end{flowtext}

\chapter{Reference}

\section{Patch parameter overview}\label{patchref}

\footnotesize
\renewcommand{\arraystretch}{1.3}
\begin{tabular}{ p{1.5cm}|p{3cm}|p{5cm}|p{5cm}|p{4cm}} 
   Number & Group & \makebox{1st press} & \makebox{2nd press} & \makebox{3rd press}\\ \hline
  0 & Speed & \makebox{Seq/Arp Speed} \linebreak \textit{numeric 0...99} & &  \\ \hline
  1 & LFO & \makebox{LFO Shape} \linebreak \textit{tri-pulse, sin-random, saw-noise} & \makebox{LFO Target} \linebreak \textit{A\&B, A, B, A \& B \& VCA } &  \makebox{LFO Clock Sync} \linebreak \textit{off, key, 1, 2, 3, 4, 5, 6, 8} \\ \hline
  2 & \makebox{Vibrato} & \makebox{Vibrato Speed} \linebreak \textit{numeric 0...99} & \makebox{Vibrato Amount} \linebreak \textit{numeric 0...99} & \makebox{Vibrato Target} \linebreak \textit{VCO, VCA, VCO A, VCO B} \\   \hline
  3 & \makebox{Modulation Wheel} & \makebox{Modulation Wheeel Target} \linebreak \textit{LFO, vibrato} & \makebox{Modulation Delay} \linebreak \textit{numeric 0...99} & \makebox{Moduation Wheel Range} \linebreak \textit{touch, soft, high, full} \\ \hline
  4 & \makebox{Configuration} & \makebox{OSC Pitch Mode} \linebreak \textit{free, semi-tones, octaves} & \makebox{External Voltage} \linebreak \textit{numeric 0...99} & \makebox{Pulse Reset Bug} \linebreak \textit{on, off} \\ \hline
  5 & Envelope & \makebox{Filter Envelope Shape} \linebreak \textit{lin-slow, lin-fast, lin-slo, exp-fast}  & \makebox{2nd Envelope Shape} \linebreak \textit{fast-lin, fast-exp, slo-lin, slo-exp} &
  \makebox{Envelope Routing} \linebreak \textit{std, poly-amp, poly, gate}\\ \hline
  6 & Bender & \makebox{Bender Target} \linebreak \textit{off, AB, VCO, VCF, vol, B} & \makebox{Bend Range}  \linebreak \textit{2nd, 3rd, 5th, octave} &  \\ \hline
  7 & \makebox{Play Mode} & \makebox{Note Priority} \linebreak \textit{last, low, high} & \makebox{Voice Assigner} \linebreak \textit{first, cycle} & Glide\textsuperscript{*)} \linebreak \textit{numeric 0...99}  \\ \hline
  8 & \makebox{Oscillator} & Detune \linebreak \textit{numeric 0...99}  &   \makebox{Vintage (Spread)} \linebreak \textit{numeric 0...99} & Drive\textsuperscript{*)} \linebreak \textit{numeric 0...99} \\ \hline
  9 & \makebox{Touch Sensitivity} & \makebox{Amplitude Velocity} \linebreak \textit{numeric 0...99}  & \makebox{Filter Velocity} \linebreak \textit{numeric 0...99} &  \\
  
\end{tabular}

\normalsize

\textsuperscript{*)} The availability of parameters \glide and \drive depend on the panel layout. In \textit{GliGli} panel layout \glide is available and drive is obsolete. In \textit{SCI} layout the parameter \drive is available while the glide parameter is set using the dedicated dial on the panel.  


\section{Settings parameter overview}\label{settingsref}

\footnotesize
\renewcommand{\arraystretch}{1.6}
\begin{tabular}{ p{2cm}|p{3cm}|p{8cm}|p{8cm}} 
   Number & Purpose & \makebox{1st press} & \makebox{Repeated press} \\
 \hline
  0 & \makebox{MIDI mode} & \makebox{Show current selected voice status} & \makebox{Toggle MIDI mode} \linebreak \textit{local on, local off}  \\
 \hline
  1 & \makebox{MIDI Receive} & \makebox{Show current MIDI receive channel} & \makebox{Cycle through channels} \linebreak \textit{OMNI, Ch1 ... Ch16} \\
 \hline
  2 & \makebox{MIDI Send} & \makebox{Show current MIDI send channel} & \makebox{Cycle through channels} \linebreak \textit{Ch1 ... Ch16} \\
 \hline
  3 & \makebox{Bender Calibration} & \makebox{Activate calibration mode} & \makebox{Confirmation} \\
 \hline
  4 & \makebox{Voice Selection} & \makebox{Show current selected voice} & \makebox{Cycle through voices} \linebreak \textit{1 ... 6} \\
 \hline
  5 & \makebox{Voice Deactivate} & \makebox{Show current selected voice status} & \makebox{Toggle voice status} \linebreak \textit{on, off} \\
 \hline
  6 & \makebox{Panel Layout} & \makebox{Show current selected panel layout} & \makebox{Toggle panel layout} \linebreak \textit{GliGli, SCI} \\
 \hline
  7 & \makebox{Unused} & & \\
 \hline
  8 & \makebox{Clock Sync} & \makebox{Show current clock sync} & \makebox{Cycle through sync types} \linebreak \textit{internal, MIDI, tape}  \\
 \hline
  9 & \makebox{Spread \& VCF limit} & \makebox{Show current setting} & \makebox{Cycle through settings} \linebreak \textit{on-on, on-off, off-on, off-off}  \\
 \hline
  Preset & \makebox{Default Patch} & \makebox{Activate default patch load} & \makebox{Confirmation} \\
 \hline
  Tune & \makebox{Per Note Tuning} & \makebox{Activate per note tuning} & \makebox{Deactivate/save per note tuning} \\
 \hline
  Record & \makebox{MIDI Patch Receive Mode} & \makebox{Activate MIDI patch receive mode} & \makebox{End MIDI patch receive mode} \\
  \end{tabular}
\normalsize


\section{MIDI control implementation}\label{midiimplenentation}

Two types of MIDI events are implemented, e.g. the upgraded Prophet 600 can receive and apply them:

\begin{itemize}
  \setlength\itemsep{0cm}
  \item MIDI CC Continuous parameters: value resolution of 0-16383 using 2 CCs (fine and coarse), or value resolution of 0-127 using only the coarse one
  \item MIDI CC Stepped parameters: 0-127, variable number of steps. They work by dividing the 0-127 range in as many zones as there are choices for the parameter. E.g.: "Unison" is an on-off parameter and therefore has s choices. In this case it is \textit{off} for 0-63 and it is \textbf{on} for 64-127.
  \item Performance events (note events, pitch bend)
  \item Real time events
  \item SysEx (for different purposes)
\end{itemize}

The Prophet 600 receives Continuous Controllers (CC) in \presetmode only. 

\textbf{MIDI CC Patch Parameters} 

\footnotesize
\renewcommand{\arraystretch}{1.3}

\begin{longtable}[l]{ p{5cm}|p{2cm}|p{1.5cm}|p{1.5cm}|p{5cm}|p{2cm}|p{1cm}} 
\textbf{Continuous Parameter} & \textbf{Type} & \textbf{CC Coarse} & \textbf{CC Fine} & \textbf{Stepped Parameter} & \textbf{Type} & \textbf{CC} \\ \hline
\endfirsthead
\textbf{Continuous Parameter} & \textbf{Type} & \textbf{CC Coarse} & \textbf{CC Fine} & \textbf{Stepped Parameter} & \textbf{Type} & \textbf{CC} \\ \hline
\endhead 
Osc A Frequency & Continuous & 16 & 80 & Osc A Saw & Stepped & 48 \\ \hline
Osc A Volume & Continuous & 17 & 81 & Osc A Triangle & Stepped & 49 \\ \hline
Osc A Pulse Width & Continuous & 18 & 82 & Osc A Square & Stepped & 50 \\ \hline
Osc B Frequency & Continuous & 19 & 83 & Osc B Saw & Stepped & 51 \\ \hline
Osc B Volume & Continuous & 20 & 84 & Osc B Triangle & Stepped & 52 \\ \hline
Osc B Pulse Width & Continuous & 21 & 85 & Osc B Sqr & Stepped & 53 \\ \hline
Osc B Fine & Continuous & 22 & 86 & Sync & Stepped & 54 \\ \hline
Cutoff & Continuous & 23 & 87 & Poly Mod Oscillator A Destination & Stepped & 55 \\ \hline
Resonance & Continuous & 24 & 88 & Poly Mod Filter Destination & Stepped & 56 \\ \hline
Filter Envelope Amount & Continuous & 25 & 89 & LFO Shape & Stepped & 57 \\ \hline
Filter Release  &  Continuous  & 26 & 90 &  LFO Targets  &  Stepped  &  59 \\ \hline
Filter Sustain  &  Continuous  & 27 & 91 &  Keyboard Filter Tracking  &  Stepped  &  60 \\ \hline
Filter Decay  &  Continuous  & 28 & 92 &  Filter Envelope Shape  &  Stepped  &  61 \\ \hline
Filter Attack  &  Continuous  & 29 & 93 &  Filter Envelope Fast/Slow  &  Stepped  &  62 \\ \hline
2nd Release  &  Continuous  & 30 & 94 &  2nd Envelope Shape  &  Stepped  &  63 \\ \hline
2nd Sustain  &  Continuous  & 31 & 95 &  Unison  &  Stepped  &  65 \\ \hline
2nd Decay  &  Continuous  & 32 & 96 &  Assigner Priority Mode  &  Stepped  &  66 \\ \hline
2nd Attack  &  Continuous  & 33 & 97 &  Bender Range (semitones)  &  Stepped  &  67 \\ \hline
Poly Mod Filter Amount  &  Continuous  & 34 & 98 &  Bender Target  &  Stepped  &  68 \\ \hline
Poly Mod Osc B Amount  &  Continuous  & 35 & 99 &  Mod Wheel Range  &  Stepped  &  69 \\ \hline
LFO Frequency  &  Continuous  & 36 & 100 &  Osc pitch mode  &  Stepped  &  70 \\ \hline
LFO Amount  &  Continuous  & 37 & 101 &  Mod Wheel Target  &  Stepped  &  71 \\ \hline
Glide  &  Continuous  & 38 & 102 &  Vibrato Target  &  Stepped  &  72 \\ \hline
Amp Velocity  &  Continuous  & 39 & 103 &   2nd Envelope Fast/Slow  &  Stepped  &  73 \\ \hline
Filter Velocity  &  Continuous  & 40 & 104 &  Sync Bug  &  Stepped  &  74 \\ \hline
Modulation delay  &  Continuous  & 41 & 105 &  Voice Assign  &  Stepped  &  75 \\ \hline
Vibrato frequency  &  Continuous  & 42 & 106 &  Envelope Routing  &  Stepped  &  76 \\ \hline
Vibrato amount  &  Continuous  & 43 & 107 &  LFO Sync  &  Stepped  &  77 \\ \hline
Unison detune  &  Continuous  & 44 & 108 &    &    &   \\ \hline
External CV amount  &  Continuous  & 46 & 110 &  &  &   \\ \hline 
Vintage / Spread  &  Continuous  & 47 & 111 &  &  &   \\ \hline 
 
\end{longtable}

\normalsize

\textbf{Change in MIDI CC from version 2.1 RC3 to \version}

The following MIDI CC have been changed.

\begin{itemize}
  \item For the filter envelope release and decay and the 2nd envelope release and decay the time value range has been rescaled for the \textit{linear} envelope shape. This was an unavoidable side effect of changing the strictly linear shapes to linear shapes with soft tails. The exponential envelopes are unaffected.
  \item The \lfofreq has been totally rescaled. The stepped frequency range parameter has been omitted and the value range of the \lfofreq dial now covers the entire frequency range. However, the dial to frequency relation has also been changed from linear to exponential, making it much softer at low settings. Consequently, the MIDI CC support for the LFO frequency range (up to firmware version 2.1 RC3) has been omitted. In order to avoid potential conflicts with MIDI tools developed for firmware versions 2.0 and 2.1 RC3, the legacy MIDI CC 58 has not been re-designated. MIDI CC 58 is ignored.
  \item The value to frequency relation for the \vibspeed has been changed from linear to exponential.   
  \item The \modwheelrange has been reconfigured with new value which also produce different wheel actions.
\end{itemize}

\textbf{Other MIDI CC Settings}

\begin{longtable}[l]{ p{6cm}|p{2.5cm}|p{1.5cm}|p{2cm}} 
\textbf{Parameter} & \textbf{Type} & \textbf{CC} & \textbf{CC Fine} \\ \hline
Toggle Live/Preset Mode & Stepped & 0 & \\ \hline
Mod Wheel & Continuous & 1 & \\ \hline
Master Volume & Continuous & 7 & \\ \hline
Sequencer / arpeggiator clock & Continuous & 45 & 109 \\ \hline
Hold Pedal & Stepped & 64 & \\ \hline
Local On/Off & Stepped & 120 & \\ \hline
All Notes Off & Stepped & 123 & \\ \hline
\end{longtable}



\normalsize


\chapter{Technical reference for the GliGli upgrade}

\section{Clock cycle}

(this nees to be done)

\section{Settings storage}

(this nees to be done)

\section{Patch storage}

(this nees to be done)

\chapter{Troubleshooting}

(this nees to be done)

\begin{thebibliography}{50} 

\bibitem{gligli}GliGli Projects Page, http://gliglisynth.blogspot.com/
\bibitem{curtis}Curtis Electro Music on Webarchive, https://web.archive.org/web/20130502132536/http://curtiselectromusic.com/Customers\_and\_Instruments.html
\bibitem{versiontwo}GliGli upgrade release 2.0,  https://gligli.github.io/p600fw/
\bibitem{modinstructions} Prophet600 Modification Instructions
\bibitem{projectsite} Project Site: http://gligli.github.io/p600fw/
\bibitem{repository} Repository: https://github.com/gligli/p600fw/
\bibitem{teensy}Teensy Site: http://www.pjrc.com/teensy/index.html
\bibitem{p600siownersmanual}Prophet 600 Owner's Manual, Sequential Circuits
\bibitem{p600siservicemanual} Prophet 600 Service Manual, Sequential Circuits

\end{thebibliography}

\end{document}
